% ============================================================
% Question Bank: ch01-02 Recognizing Proportional Relationships
% Topic Title: Recognizing Proportional Relationships
% CCSS: 7.RP.A.2.A
% Questions: 30 (18 MC + 7 SA + 5 Graphical)
% ============================================================

% --- Multiple Choice Questions (18) ---

\begin{practiceQuestion}{1-2-q01}{mc}
\begin{questionText}
Which table shows a proportional relationship?
\end{questionText}
\choiceA{\begin{tabular}{|c|c|} \hline $x$ & $y$ \\ \hline $2$ & $5$ \\ $4$ & $10$ \\ $6$ & $15$ \\ \hline \end{tabular}}
\choiceB{\begin{tabular}{|c|c|} \hline $x$ & $y$ \\ \hline $1$ & $3$ \\ $2$ & $5$ \\ $3$ & $7$ \\ \hline \end{tabular}}
\choiceC{\begin{tabular}{|c|c|} \hline $x$ & $y$ \\ \hline $2$ & $6$ \\ $3$ & $8$ \\ $4$ & $10$ \\ \hline \end{tabular}}
\choiceD{\begin{tabular}{|c|c|} \hline $x$ & $y$ \\ \hline $1$ & $4$ \\ $2$ & $6$ \\ $3$ & $8$ \\ \hline \end{tabular}}
\correctAnswer{A}
\explanation{In choice A, $\frac{y}{x} = \frac{5}{2} = 2.5$ for every row. The other tables have ratios that change from row to row.}
\end{practiceQuestion}

\begin{practiceQuestion}{1-2-q02}{mc}
\begin{questionText}
A graph of a proportional relationship must:
\end{questionText}
\choiceA{Be a straight line}
\choiceB{Pass through the origin $(0, 0)$}
\choiceC{Be a straight line that passes through $(0, 0)$}
\choiceD{Have a positive slope}
\correctAnswer{C}
\explanation{A proportional relationship graphs as a straight line AND it must pass through the origin. Both conditions are needed.}
\end{practiceQuestion}

\begin{practiceQuestion}{1-2-q03}{mc}
\begin{questionText}
Sara earns $\$9$ per hour with no signing bonus. Tom earns $\$8$ per hour plus a $\$20$ signing bonus. Which pay plan is proportional?
\end{questionText}
\choiceA{Sara's}
\choiceB{Tom's}
\choiceC{Both}
\choiceD{Neither}
\correctAnswer{A}
\explanation{Sara's pay is proportional because her total pay equals a constant rate times hours ($p = 9h$). Tom's has a fixed $\$20$ bonus, which means $0$ hours still gives $\$20$, so his is not proportional.}
\end{practiceQuestion}

\begin{practiceQuestion}{1-2-q04}{mc}
\begin{questionText}
The table shows gallons of gas used and miles driven. Is this proportional?

\begin{center}
\begin{tabular}{|c|c|}
\hline \textbf{Gallons} & \textbf{Miles} \\ \hline
$2$ & $56$ \\ $5$ & $140$ \\ $8$ & $224$ \\ \hline
\end{tabular}
\end{center}
\end{questionText}
\choiceA{Yes, because the differences are equal}
\choiceB{Yes, because $\frac{y}{x} = 28$ for every row}
\choiceC{No, because the values increase}
\choiceD{No, because $8 - 5 \neq 5 - 2$}
\correctAnswer{B}
\explanation{$56 \div 2 = 28$, $140 \div 5 = 28$, $224 \div 8 = 28$. Every ratio is $28$, so this is proportional.}
\end{practiceQuestion}

\begin{practiceQuestion}{1-2-q05}{mc}
\begin{questionText}
Which situation does NOT represent a proportional relationship?
\end{questionText}
\choiceA{Buying apples at $\$1.50$ each}
\choiceB{A taxi that charges $\$3$ base fee plus $\$2$ per mile}
\choiceC{Earning $\$12$ per hour}
\choiceD{Printing $25$ pages per minute}
\correctAnswer{B}
\explanation{The taxi has a $\$3$ base fee. At $0$ miles, the cost is $\$3$ (not $\$0$), so $\frac{y}{x}$ is not constant and the relationship is not proportional.}
\end{practiceQuestion}

\begin{practiceQuestion}{1-2-q06}{mc}
\begin{questionText}
A line passes through $(0, 3)$ and $(2, 7)$. Is the relationship proportional?
\end{questionText}
\choiceA{Yes, because it is a straight line}
\choiceB{Yes, because the slope is constant}
\choiceC{No, because it does not pass through the origin}
\choiceD{No, because the slope is too steep}
\correctAnswer{C}
\explanation{Although the graph is a straight line, it passes through $(0, 3)$ instead of $(0, 0)$. A proportional relationship must pass through the origin.}
\end{practiceQuestion}

\begin{practiceQuestion}{1-2-q07}{mc}
\begin{questionText}
Is the following relationship proportional?

\begin{center}
\begin{tabular}{|c|c|}
\hline $x$ & $y$ \\ \hline $4$ & $12$ \\ $6$ & $18$ \\ $9$ & $27$ \\ \hline
\end{tabular}
\end{center}
\end{questionText}
\choiceA{No, because $y$ is always larger than $x$}
\choiceB{No, because the differences between rows are not equal}
\choiceC{Yes, because $\frac{y}{x} = 3$ for every row}
\choiceD{Yes, because $y$ increases as $x$ increases}
\correctAnswer{C}
\explanation{$12 \div 4 = 3$, $18 \div 6 = 3$, $27 \div 9 = 3$. The ratio is constant at $3$, so the relationship is proportional.}
\end{practiceQuestion}

\begin{practiceQuestion}{1-2-q08}{mc}
\begin{questionText}
A student says the table below is proportional because the $y$-values go up by $4$ each time. What is wrong with this reasoning?

\begin{center}
\begin{tabular}{|c|c|}
\hline $x$ & $y$ \\ \hline $1$ & $5$ \\ $2$ & $9$ \\ $3$ & $13$ \\ \hline
\end{tabular}
\end{center}
\end{questionText}
\choiceA{The values should go up by the same amount}
\choiceB{A constant difference does not mean a constant ratio}
\choiceC{The table is actually proportional}
\choiceD{You cannot use tables to test proportionality}
\correctAnswer{B}
\explanation{The ratios are $5 \div 1 = 5$, $9 \div 2 = 4.5$, $13 \div 3 \approx 4.3$. They are not equal. A constant difference (additive change) does not guarantee a constant ratio (multiplicative change).}
\end{practiceQuestion}

\begin{practiceQuestion}{1-2-q09}{mc}
\begin{questionText}
Which test determines if a relationship is proportional using a table?
\end{questionText}
\choiceA{Check that $y$ increases by the same amount each time}
\choiceB{Check that $\frac{y}{x}$ is the same for every row}
\choiceC{Check that $x + y$ is constant}
\choiceD{Check that $y - x$ is constant}
\correctAnswer{B}
\explanation{For a proportional relationship, the ratio $\frac{y}{x}$ must be the same for every pair. This is the table test for proportionality.}
\end{practiceQuestion}

\begin{practiceQuestion}{1-2-q10}{mc}
\begin{questionText}
A gym charges $\$25$ per month with no enrollment fee. Is the relationship between months and total cost proportional?
\end{questionText}
\choiceA{No, because $\$25$ is too much}
\choiceB{No, because the cost goes up each month}
\choiceC{Yes, because total cost $= 25 \times$ months with no extra fee}
\choiceD{Yes, because it is a real-world situation}
\correctAnswer{C}
\explanation{With no enrollment fee, at $0$ months the cost is $\$0$. Total cost $= 25m$, making $\frac{y}{x} = 25$ for every month. This is proportional.}
\end{practiceQuestion}

\begin{practiceQuestion}{1-2-q11}{mc}
\begin{questionText}
Is this relationship proportional?

\begin{center}
\begin{tabular}{|c|c|}
\hline $x$ & $y$ \\ \hline $0$ & $0$ \\ $3$ & $9$ \\ $5$ & $16$ \\ \hline
\end{tabular}
\end{center}
\end{questionText}
\choiceA{Yes, because it starts at $(0, 0)$}
\choiceB{Yes, because $y$ is always bigger than $x$}
\choiceC{No, because $\frac{9}{3} = 3$ but $\frac{16}{5} = 3.2$}
\choiceD{No, because $x$ and $y$ are different}
\correctAnswer{C}
\explanation{The ratio $\frac{y}{x}$ changes: $\frac{9}{3} = 3$ and $\frac{16}{5} = 3.2$. Since the ratios are not equal, this is not proportional — even though it passes through $(0, 0)$.}
\end{practiceQuestion}

\begin{practiceQuestion}{1-2-q12}{mc}
\begin{questionText}
The number of pages Tyler reads is proportional to the time he spends reading. If he reads $15$ pages in $20$ minutes, which of the following must also be true?
\end{questionText}
\choiceA{He reads $25$ pages in $30$ minutes}
\choiceB{He reads $30$ pages in $40$ minutes}
\choiceC{He reads $20$ pages in $25$ minutes}
\choiceD{He reads $10$ pages in $10$ minutes}
\correctAnswer{B}
\explanation{The unit rate is $\frac{15}{20} = 0.75$ pages/min. Check: $0.75 \times 40 = 30$ pages. Only choice B keeps the same ratio $\frac{y}{x} = 0.75$.}
\end{practiceQuestion}

\begin{practiceQuestion}{1-2-q13}{mc}
\begin{questionText}
A pool fills at a constant rate. Which pair of conditions proves the relationship between time and water volume is proportional?
\end{questionText}
\choiceA{The water level rises by the same amount each hour}
\choiceB{The pool starts empty and the rate is constant}
\choiceC{The pool starts with some water and the rate is constant}
\choiceD{The graph is a curve through the origin}
\correctAnswer{B}
\explanation{A proportional relationship requires $y = kx$, meaning it starts at $(0, 0)$ (pool starts empty) and has a constant rate. Option C fails because starting with water gives a non-zero $y$-intercept.}
\end{practiceQuestion}

\begin{practiceQuestion}{1-2-q14}{mc}
\begin{questionText}
Which of these graphs shows a proportional relationship?
\end{questionText}
\choiceA{A straight line through $(0, 2)$ and $(3, 8)$}
\choiceB{A straight line through $(0, 0)$ and $(4, 12)$}
\choiceC{A curve through $(0, 0)$ and $(4, 16)$}
\choiceD{A horizontal line through $(0, 5)$}
\correctAnswer{B}
\explanation{A proportional relationship is a straight line through the origin. Only B satisfies both conditions. A doesn't start at the origin, C is a curve, and D is horizontal (zero rate is technically proportional but passes through $(0, 5)$, not through the origin in the usual sense).}
\end{practiceQuestion}

\begin{practiceQuestion}{1-2-q15}{mc}
\begin{questionText}
A table shows that $\frac{y}{x} = 7$ for the first two rows, but $\frac{y}{x} = 7.5$ for the third row. What can you conclude?
\end{questionText}
\choiceA{The relationship is proportional with $k = 7$}
\choiceB{The relationship is proportional with $k = 7.5$}
\choiceC{The relationship is not proportional}
\choiceD{You need more rows to decide}
\correctAnswer{C}
\explanation{Every ratio must be the same for a proportional relationship. One mismatch is enough to prove it is not proportional.}
\end{practiceQuestion}

\begin{practiceQuestion}{1-2-q16}{mc}
\begin{questionText}
Maria says the table below shows a proportional relationship because $y$ is always $3$ more than $x$. Is she correct?

\begin{center}
\begin{tabular}{|c|c|}
\hline $x$ & $y$ \\ \hline $2$ & $5$ \\ $4$ & $7$ \\ $6$ & $9$ \\ \hline
\end{tabular}
\end{center}
\end{questionText}
\choiceA{Yes, because $y - x = 3$ always}
\choiceB{Yes, because $y$ increases as $x$ increases}
\choiceC{No, because $\frac{5}{2} \neq \frac{7}{4} \neq \frac{9}{6}$}
\choiceD{No, because $x$ and $y$ are both even or odd}
\correctAnswer{C}
\explanation{$\frac{5}{2} = 2.5$, $\frac{7}{4} = 1.75$, $\frac{9}{6} = 1.5$. The ratios are not equal. A constant difference ($y = x + 3$) does not make a proportional relationship.}
\end{practiceQuestion}

\begin{practiceQuestion}{1-2-q17}{mc}
\begin{questionText}
A phone plan charges $\$0.10$ per text with no monthly fee. Another plan charges $\$5$ per month plus $\$0.05$ per text. Which statement is true?
\end{questionText}
\choiceA{Both plans are proportional}
\choiceB{Only the first plan is proportional}
\choiceC{Only the second plan is proportional}
\choiceD{Neither plan is proportional}
\correctAnswer{B}
\explanation{The first plan has no fixed fee: cost $= 0.10t$, which is proportional. The second plan has a $\$5$ monthly fee that doesn't depend on texts, making it non-proportional.}
\end{practiceQuestion}

\begin{practiceQuestion}{1-2-q18}{mc}
\begin{questionText}
Is the following relationship proportional?

\begin{center}
\begin{tabular}{|c|c|}
\hline $x$ & $y$ \\ \hline $5$ & $8$ \\ $10$ & $16$ \\ $15$ & $24$ \\ \hline
\end{tabular}
\end{center}
\end{questionText}
\choiceA{No, because $8 - 5 \neq 16 - 10$}
\choiceB{No, because the $x$-values skip by $5$}
\choiceC{Yes, because $\frac{y}{x} = 1.6$ for every row}
\choiceD{Yes, because $y$ doubles when $x$ doubles}
\correctAnswer{C}
\explanation{$\frac{8}{5} = 1.6$, $\frac{16}{10} = 1.6$, $\frac{24}{15} = 1.6$. Every ratio is the same, so the relationship is proportional.}
\end{practiceQuestion}

% --- Short Answer Questions (7) ---

\begin{practiceQuestion}{1-2-q19}{sa}
\begin{questionText}
Is the relationship in the table below proportional? Explain with the ratios.

\begin{center}
\begin{tabular}{|c|c|}
\hline $x$ & $y$ \\ \hline $3$ & $9$ \\ $7$ & $21$ \\ $10$ & $30$ \\ \hline
\end{tabular}
\end{center}
\end{questionText}
\correctAnswer{Yes, $\frac{y}{x} = 3$ for every row}
\explanation{$\frac{9}{3} = 3$, $\frac{21}{7} = 3$, $\frac{30}{10} = 3$. The ratio is the same in every row, so the relationship is proportional.}
\end{practiceQuestion}

\begin{practiceQuestion}{1-2-q20}{sa}
\begin{questionText}
Is the relationship in the table below proportional? Explain.

\begin{center}
\begin{tabular}{|c|c|}
\hline $x$ & $y$ \\ \hline $2$ & $8$ \\ $3$ & $11$ \\ $5$ & $17$ \\ \hline
\end{tabular}
\end{center}
\end{questionText}
\correctAnswer{No, the ratios are not equal}
\explanation{$\frac{8}{2} = 4$, $\frac{11}{3} \approx 3.67$, $\frac{17}{5} = 3.4$. The ratios differ, so this is not proportional.}
\end{practiceQuestion}

\begin{practiceQuestion}{1-2-q21}{sa}
\begin{questionText}
A straight-line graph passes through $(0, 0)$ and $(6, 15)$. Is this a proportional relationship? If so, what is the constant ratio?
\end{questionText}
\correctAnswer{Yes, the constant ratio is $2.5$}
\explanation{The graph is a straight line through the origin, so it is proportional. The constant ratio is $\frac{15}{6} = 2.5$.}
\end{practiceQuestion}

\begin{practiceQuestion}{1-2-q22}{sa}
\begin{questionText}
A movie streaming service charges $\$8$ per month plus $\$2$ per movie. Is the total cost proportional to the number of movies? Explain.
\end{questionText}
\correctAnswer{No, there is a fixed $\$8$ monthly fee}
\explanation{At $0$ movies, the cost is $\$8$ (not $\$0$). A proportional relationship must start at zero. The fixed fee breaks the constant ratio.}
\end{practiceQuestion}

\begin{practiceQuestion}{1-2-q23}{sa}
\begin{questionText}
Give an example of two quantities that have a proportional relationship and explain why.
\end{questionText}
\correctAnswer{Answers vary; e.g., total cost and number of items at a fixed price per item}
\explanation{Any situation where $y = kx$ with no added constant is proportional. For example, buying pencils at $\$0.50$ each: $0$ pencils cost $\$0$, and the ratio cost/pencils is always $0.50$.}
\end{practiceQuestion}

\begin{practiceQuestion}{1-2-q24}{sa}
\begin{questionText}
A table has these values: $(2, 10)$, $(4, 20)$, $(6, 30)$. Could a fifth data point of $(8, 42)$ be part of this proportional relationship? Explain.
\end{questionText}
\correctAnswer{No, because $\frac{42}{8} = 5.25 \neq 5$}
\explanation{The first three points have $\frac{y}{x} = 5$. The point $(8, 42)$ gives $\frac{42}{8} = 5.25$, which is different. To stay proportional, the value at $x = 8$ must be $40$.}
\end{practiceQuestion}

\begin{practiceQuestion}{1-2-q25}{sa}
\begin{questionText}
A straight line passes through $(0, 4)$ and $(2, 12)$. Is this proportional? Why or why not?
\end{questionText}
\correctAnswer{No, because the line does not pass through $(0, 0)$}
\explanation{Even though the graph is a straight line, it crosses the $y$-axis at $4$, not $0$. A proportional relationship must go through the origin.}
\end{practiceQuestion}

% --- Graphical Questions (3 GMC + 2 GSA) ---

\begin{practiceQuestion}{1-2-q26}{gmc}
\begin{questionText}
Look at the graph below. Is this a proportional relationship?

\begin{center}
\begin{tikzpicture}[scale=0.8]
  \draw[step=1,gray!20,thin] (0,0) grid (6,6);
  \draw[thick,->] (0,0) -- (6.5,0) node[right]{\small $x$};
  \draw[thick,->] (0,0) -- (0,6.5) node[above]{\small $y$};
  \foreach \x in {1,...,6} \draw (\x,0.07) -- (\x,-0.07) node[below]{\tiny \x};
  \foreach \y in {1,...,6} \draw (0.07,\y) -- (-0.07,\y) node[left]{\tiny \y};
  \draw[thick,funBlue] (0,2) -- (5,6);
  \fill[funBlue] (0,2) circle (2.5pt);
  \fill[funBlue] (5,6) circle (2.5pt);
\end{tikzpicture}
\end{center}
\end{questionText}
\choiceA{Yes, because it is a straight line}
\choiceB{Yes, because it goes up to the right}
\choiceC{No, because it does not pass through $(0, 0)$}
\choiceD{No, because the slope is not $1$}
\correctAnswer{C}
\explanation{The line crosses the $y$-axis at $(0, 2)$, not $(0, 0)$. A proportional graph must be a straight line through the origin.}
\end{practiceQuestion}

\begin{practiceQuestion}{1-2-q27}{gmc}
\begin{questionText}
Which graph represents a proportional relationship?

\begin{center}
\begin{tikzpicture}[scale=0.55]
  % Graph I
  \begin{scope}[xshift=0cm]
    \draw[step=1,gray!20,thin] (0,0) grid (4,4);
    \draw[thick,->] (0,0) -- (4.5,0);
    \draw[thick,->] (0,0) -- (0,4.5);
    \draw[thick,funBlue] (0,0) -- (4,3);
    \fill[funBlue] (0,0) circle (3pt); \fill[funBlue] (4,3) circle (3pt);
    \node[below] at (2,-0.5) {\small I};
  \end{scope}
  % Graph II
  \begin{scope}[xshift=6cm]
    \draw[step=1,gray!20,thin] (0,0) grid (4,4);
    \draw[thick,->] (0,0) -- (4.5,0);
    \draw[thick,->] (0,0) -- (0,4.5);
    \draw[thick,funRed] (0,1) -- (4,4);
    \fill[funRed] (0,1) circle (3pt); \fill[funRed] (4,4) circle (3pt);
    \node[below] at (2,-0.5) {\small II};
  \end{scope}
  % Graph III
  \begin{scope}[xshift=12cm]
    \draw[step=1,gray!20,thin] (0,0) grid (4,4);
    \draw[thick,->] (0,0) -- (4.5,0);
    \draw[thick,->] (0,0) -- (0,4.5);
    \draw[thick,funGreen] (0,0) parabola (4,4);
    \fill[funGreen] (0,0) circle (3pt); \fill[funGreen] (4,4) circle (3pt);
    \node[below] at (2,-0.5) {\small III};
  \end{scope}
\end{tikzpicture}
\end{center}
\end{questionText}
\choiceA{Graph I only}
\choiceB{Graph II only}
\choiceC{Graph III only}
\choiceD{Graphs I and III}
\correctAnswer{A}
\explanation{Graph I is a straight line through the origin — proportional. Graph II is a straight line but starts at $(0, 1)$. Graph III passes through the origin but is curved, not straight.}
\end{practiceQuestion}

\begin{practiceQuestion}{1-2-q28}{gmc}
\begin{questionText}
The table below shows the cost of renting kayaks. Is this a proportional relationship?

\begin{center}
\begin{tabular}{|c|c|c|}
\hline
\textbf{Hours ($x$)} & \textbf{Cost ($y$)} & $\frac{y}{x}$ \\ \hline
$1$ & $\$12$ & $12$ \\ \hline
$2$ & $\$24$ & $12$ \\ \hline
$3$ & $\$36$ & $12$ \\ \hline
$4$ & $\$48$ & $12$ \\ \hline
\end{tabular}
\end{center}
\end{questionText}
\choiceA{No, because the cost keeps changing}
\choiceB{No, because $\$12$ is not a whole-number ratio}
\choiceC{Yes, because $\frac{y}{x} = 12$ for every row}
\choiceD{Yes, because the cost goes up by $\$12$ each time}
\correctAnswer{C}
\explanation{The ratio $\frac{y}{x} = 12$ for every row. A constant ratio means proportional. Note that a constant difference (D) is not sufficient by itself.}
\end{practiceQuestion}

\begin{practiceQuestion}{1-2-q29}{gsa}
\begin{questionText}
Look at the graph below. Is this relationship proportional? Explain your reasoning.

\begin{center}
\begin{tikzpicture}[scale=0.8]
  \draw[step=1,gray!20,thin] (0,0) grid (6,6);
  \draw[thick,->] (0,0) -- (6.5,0) node[right]{\small $x$};
  \draw[thick,->] (0,0) -- (0,6.5) node[above]{\small $y$};
  \foreach \x in {1,...,6} \draw (\x,0.07) -- (\x,-0.07) node[below]{\tiny \x};
  \foreach \y in {1,...,6} \draw (0.07,\y) -- (-0.07,\y) node[left]{\tiny \y};
  \draw[thick,funBlue] (0,0) -- (6,4);
  \fill[funBlue] (0,0) circle (2.5pt);
  \fill[funBlue] (3,2) circle (2.5pt);
  \fill[funBlue] (6,4) circle (2.5pt);
\end{tikzpicture}
\end{center}
\end{questionText}
\correctAnswer{Yes, it is a straight line through the origin}
\explanation{The graph is a straight line that passes through $(0, 0)$. Both conditions for a proportional relationship are met: straight line and passes through the origin.}
\end{practiceQuestion}

\begin{practiceQuestion}{1-2-q30}{gsa}
\begin{questionText}
The table below shows a relationship between $x$ and $y$. Determine whether the relationship is proportional. If not, change one $y$-value to make it proportional.

\begin{center}
\begin{tabular}{|c|c|}
\hline $x$ & $y$ \\ \hline
$2$ & $6$ \\ \hline
$5$ & $15$ \\ \hline
$8$ & $25$ \\ \hline
\end{tabular}
\end{center}
\end{questionText}
\correctAnswer{Not proportional; change $25$ to $24$}
\explanation{$\frac{6}{2} = 3$, $\frac{15}{5} = 3$, $\frac{25}{8} = 3.125$. The third ratio is different. To make it proportional with $k = 3$, the third $y$-value should be $3 \times 8 = 24$.}
\end{practiceQuestion}
